% Options for packages loaded elsewhere
\PassOptionsToPackage{unicode}{hyperref}
\PassOptionsToPackage{hyphens}{url}
\documentclass[
]{article}
\usepackage{xcolor}
\usepackage{amsmath,amssymb}
\setcounter{secnumdepth}{-\maxdimen} % remove section numbering
\usepackage{iftex}
\ifPDFTeX
  \usepackage[T1]{fontenc}
  \usepackage[utf8]{inputenc}
  \usepackage{textcomp} % provide euro and other symbols
\else % if luatex or xetex
  \usepackage{unicode-math} % this also loads fontspec
  \defaultfontfeatures{Scale=MatchLowercase}
  \defaultfontfeatures[\rmfamily]{Ligatures=TeX,Scale=1}
\fi
\usepackage{lmodern}
\ifPDFTeX\else
  % xetex/luatex font selection
\fi
% Use upquote if available, for straight quotes in verbatim environments
\IfFileExists{upquote.sty}{\usepackage{upquote}}{}
\IfFileExists{microtype.sty}{% use microtype if available
  \usepackage[]{microtype}
  \UseMicrotypeSet[protrusion]{basicmath} % disable protrusion for tt fonts
}{}
\makeatletter
\@ifundefined{KOMAClassName}{% if non-KOMA class
  \IfFileExists{parskip.sty}{%
    \usepackage{parskip}
  }{% else
    \setlength{\parindent}{0pt}
    \setlength{\parskip}{6pt plus 2pt minus 1pt}}
}{% if KOMA class
  \KOMAoptions{parskip=half}}
\makeatother
\usepackage{color}
\usepackage{fancyvrb}
\newcommand{\VerbBar}{|}
\newcommand{\VERB}{\Verb[commandchars=\\\{\}]}
\DefineVerbatimEnvironment{Highlighting}{Verbatim}{commandchars=\\\{\}}
% Add ',fontsize=\small' for more characters per line
\newenvironment{Shaded}{}{}
\newcommand{\AlertTok}[1]{\textcolor[rgb]{1.00,0.00,0.00}{\textbf{#1}}}
\newcommand{\AnnotationTok}[1]{\textcolor[rgb]{0.38,0.63,0.69}{\textbf{\textit{#1}}}}
\newcommand{\AttributeTok}[1]{\textcolor[rgb]{0.49,0.56,0.16}{#1}}
\newcommand{\BaseNTok}[1]{\textcolor[rgb]{0.25,0.63,0.44}{#1}}
\newcommand{\BuiltInTok}[1]{\textcolor[rgb]{0.00,0.50,0.00}{#1}}
\newcommand{\CharTok}[1]{\textcolor[rgb]{0.25,0.44,0.63}{#1}}
\newcommand{\CommentTok}[1]{\textcolor[rgb]{0.38,0.63,0.69}{\textit{#1}}}
\newcommand{\CommentVarTok}[1]{\textcolor[rgb]{0.38,0.63,0.69}{\textbf{\textit{#1}}}}
\newcommand{\ConstantTok}[1]{\textcolor[rgb]{0.53,0.00,0.00}{#1}}
\newcommand{\ControlFlowTok}[1]{\textcolor[rgb]{0.00,0.44,0.13}{\textbf{#1}}}
\newcommand{\DataTypeTok}[1]{\textcolor[rgb]{0.56,0.13,0.00}{#1}}
\newcommand{\DecValTok}[1]{\textcolor[rgb]{0.25,0.63,0.44}{#1}}
\newcommand{\DocumentationTok}[1]{\textcolor[rgb]{0.73,0.13,0.13}{\textit{#1}}}
\newcommand{\ErrorTok}[1]{\textcolor[rgb]{1.00,0.00,0.00}{\textbf{#1}}}
\newcommand{\ExtensionTok}[1]{#1}
\newcommand{\FloatTok}[1]{\textcolor[rgb]{0.25,0.63,0.44}{#1}}
\newcommand{\FunctionTok}[1]{\textcolor[rgb]{0.02,0.16,0.49}{#1}}
\newcommand{\ImportTok}[1]{\textcolor[rgb]{0.00,0.50,0.00}{\textbf{#1}}}
\newcommand{\InformationTok}[1]{\textcolor[rgb]{0.38,0.63,0.69}{\textbf{\textit{#1}}}}
\newcommand{\KeywordTok}[1]{\textcolor[rgb]{0.00,0.44,0.13}{\textbf{#1}}}
\newcommand{\NormalTok}[1]{#1}
\newcommand{\OperatorTok}[1]{\textcolor[rgb]{0.40,0.40,0.40}{#1}}
\newcommand{\OtherTok}[1]{\textcolor[rgb]{0.00,0.44,0.13}{#1}}
\newcommand{\PreprocessorTok}[1]{\textcolor[rgb]{0.74,0.48,0.00}{#1}}
\newcommand{\RegionMarkerTok}[1]{#1}
\newcommand{\SpecialCharTok}[1]{\textcolor[rgb]{0.25,0.44,0.63}{#1}}
\newcommand{\SpecialStringTok}[1]{\textcolor[rgb]{0.73,0.40,0.53}{#1}}
\newcommand{\StringTok}[1]{\textcolor[rgb]{0.25,0.44,0.63}{#1}}
\newcommand{\VariableTok}[1]{\textcolor[rgb]{0.10,0.09,0.49}{#1}}
\newcommand{\VerbatimStringTok}[1]{\textcolor[rgb]{0.25,0.44,0.63}{#1}}
\newcommand{\WarningTok}[1]{\textcolor[rgb]{0.38,0.63,0.69}{\textbf{\textit{#1}}}}
\usepackage{longtable,booktabs,array}
\newcounter{none} % for unnumbered tables
\usepackage{calc} % for calculating minipage widths
% Correct order of tables after \paragraph or \subparagraph
\usepackage{etoolbox}
\makeatletter
\patchcmd\longtable{\par}{\if@noskipsec\mbox{}\fi\par}{}{}
\makeatother
% Allow footnotes in longtable head/foot
\IfFileExists{footnotehyper.sty}{\usepackage{footnotehyper}}{\usepackage{footnote}}
\makesavenoteenv{longtable}
\setlength{\emergencystretch}{3em} % prevent overfull lines
\providecommand{\tightlist}{%
  \setlength{\itemsep}{0pt}\setlength{\parskip}{0pt}}
\usepackage[]{natbib}
\bibliographystyle{plainnat}
\usepackage{bookmark}
\IfFileExists{xurl.sty}{\usepackage{xurl}}{} % add URL line breaks if available
\urlstyle{same}
% fallback for those not using the hyperref driver hyperxmp:
\makeatletter
\@ifundefined{xmpquote}{\newcommand{\xmpquote}[1]{#1}}{}
\makeatother
\hypersetup{
  pdftitle={Critical-Boundary Calibration of Cosmic-Fate Probabilities},
  pdfauthor={Independent researcher},
  hidelinks,
  pdfcreator={LaTeX via pandoc}}

\title{Critical-Boundary Calibration of Cosmic-Fate Probabilities}
\author{Independent researcher}
\date{2026-07-06}

\begin{document}
\maketitle

\section{Critical-Boundary Calibration of Cosmic-Fate
Probabilities}\label{critical-boundary-calibration-of-cosmic-fate-probabilities}

\subsection{Abstract}\label{abstract}

Posterior probabilities for cosmic-fate classes are often read as direct
forecasts: one infers cosmological parameters, applies a future-fate
classifier to posterior samples, and reports the resulting class
fractions. This interpretation can be misleading near Lambda-CDM or any
other model that lies on or close to a fate-classification boundary. We
frame posterior fate probabilities as boundary statistics. In a
one-dimensional calibrated boundary model, repeated experiments produce
an approximately uniform distribution of posterior class probabilities,
with mean 0.501, variance 0.084, and KS p = 0.165 against Uniform(0,1).
A cosmology-shaped two-dimensional w0-wa analogue gives the same
diagnostic behavior, with mean 0.501 and KS p = 0.058. We then define a
finite-mock null-calibration protocol and apply it to an existing LCDM
mock case study. In that case, the direction layer is uninformative by
construction (\texttt{P\_heat\textgreater{}0.5} in 50/100 noisy mocks),
while the observed \texttt{P\_RIP=0.0018} lies below 0/100 noisy LCDM
mocks, giving a plus-one lower-tail p-value of 0.0099. We propose a
five-layer reporting standard that separates Direction, Depth, Evidence,
Boundary, and Horizon. The result is a protocol for reporting
boundary-sensitive cosmic-fate probabilities, not a direct claim about
the universe's final fate.

Keywords: cosmic fate; dark energy; boundary calibration; posterior
class probabilities; null mocks; Lambda-CDM; w0-wa; finite-mock
p-values.

\subsection{1. Introduction}\label{introduction}

Cosmic-fate statements sit at an awkward interface between data analysis
and long-future extrapolation. Modern cosmological data constrain
distances, expansion histories, and dark-energy parameters; a fate
statement then maps those inferred quantities through a classifier that
asks what happens far beyond the directly observed range. It is natural
to summarize such an analysis by a posterior class fraction, for example
\texttt{P(RIP)}, \texttt{P(DECAY)}, \texttt{P(DS)}, or a
heat-death-compatible aggregate. That number is useful, but its meaning
is not the same as an ordinary posterior probability for a smooth
parameter interval.

The reason is geometric. Fate labels are typically produced by threshold
rules: a finite-time divergence triggers one label, an asymptotic de
Sitter condition another, a collapse condition another, and an audit
channel catches cases not covered by the declared rules. A posterior
fate probability is therefore an integral of an indicator function over
a label region. Near the boundary between two such regions, a small
displacement of the posterior center can produce a large change in the
reported class mass even when the underlying parameter inference is well
calibrated.

This paper isolates that statistical problem. The object of study is not
the real universe's final state. The object is the repeated-experiment
behavior of posterior class masses when the data-generating model lies
on or close to a classification boundary. This distinction matters
because a boundary case can produce a direction rate near 50 percent
without indicating a broken inference pipeline. Conversely, a raw class
probability can still be informative if it is unusually deep in the null
distribution generated by the same pipeline.

The contribution is fourfold. First, we show analytically and by
simulation why a calibrated boundary posterior can produce a broad,
approximately uniform distribution of class probabilities. Second, we
translate the result into a cosmology-shaped w0-wa analogue and an
abstract fate-classifier framework. Third, we define a finite-mock
null-calibration rule that converts raw class probabilities into
tail-depth statements. Fourth, we propose a reporting standard that
separates Direction, Depth, Evidence, Boundary, and Horizon. This
separation is the central safeguard: it prevents a boundary-sensitive
posterior mass from being over-read as a direct fate forecast or as
model evidence.

The manuscript uses several standard cosmological reference points only
as context: CPL-like dark-energy parametrizations
\citep{ChevallierPolarski2001, Linder2003}, late-time data combinations
such as DESI BAO and supernova samples \citep{DESIDR2, Brout2022},
compressed CMB distance priors \citep{ChenHuangWang2018}, and the
long-standing literature on dark-energy fate questions
\citep{Dyson1979, KraussTurner1999, KraussStarkman2000, Caldwell2003}.
The method itself does not depend on accepting any particular current
cosmological claim.

\subsection{2. Boundary-Statistic
Formulation}\label{boundary-statistic-formulation}

Let \texttt{theta} denote the parameter vector used by a cosmological
analysis. It may include \texttt{w0}, \texttt{wa}, matter density,
curvature, nuisance parameters, or a richer description of the
dark-energy history. Let

\begin{Shaded}
\begin{Highlighting}[]
\NormalTok{C(theta) in \{CRUNCH, RIP, DS, DECAY, OTHER\}}
\end{Highlighting}
\end{Shaded}

be a deterministic classifier that maps each parameter point to one fate
label. For a label \texttt{L}, the posterior class mass is

\begin{Shaded}
\begin{Highlighting}[]
\NormalTok{q\_L(data) = P(C(theta)=L | data)}
\NormalTok{          = integral 1[C(theta)=L] p(theta | data) dtheta.}
\end{Highlighting}
\end{Shaded}

This is a posterior probability, but it is also a statistic produced by
a threshold map. Its repeated-experiment distribution can be broad when
the data-generating value \texttt{theta0} lies on a boundary between
labels.

For two labels \texttt{A} and \texttt{B}, write a local boundary in
terms of a score

\begin{Shaded}
\begin{Highlighting}[]
\NormalTok{g(theta) = 0,}
\end{Highlighting}
\end{Shaded}

with \texttt{A} on one side and \texttt{B} on the other. The
critical-boundary case is

\begin{Shaded}
\begin{Highlighting}[]
\NormalTok{g(theta0) = 0}
\end{Highlighting}
\end{Shaded}

or a near-boundary version where \texttt{g(theta0)} is small compared
with the posterior uncertainty in the boundary-normal direction.

The minimal Gaussian version makes the issue explicit. Suppose

\begin{Shaded}
\begin{Highlighting}[]
\NormalTok{x\_hat \textasciitilde{} Normal(theta0, sigma\_data)}
\NormalTok{theta | x\_hat \textasciitilde{} Normal(x\_hat, sigma\_post)}
\NormalTok{C(theta) = A if theta \textgreater{} 0, otherwise B.}
\end{Highlighting}
\end{Shaded}

The posterior class mass is

\begin{Shaded}
\begin{Highlighting}[]
\NormalTok{q\_A = P(theta \textgreater{} 0 | x\_hat) = Phi(x\_hat / sigma\_post).}
\end{Highlighting}
\end{Shaded}

If \texttt{theta0=0} and \texttt{sigma\_post=sigma\_data}, then

\begin{Shaded}
\begin{Highlighting}[]
\NormalTok{x\_hat / sigma\_data \textasciitilde{} Normal(0,1),}
\NormalTok{q\_A = Phi(Z),  Z \textasciitilde{} Normal(0,1).}
\end{Highlighting}
\end{Shaded}

By the probability integral transform \citep{Rosenblatt1952},
\texttt{q\_A\ \textasciitilde{}\ Uniform(0,1)}. Thus a perfectly
calibrated boundary null produces class probabilities spread across the
full unit interval. A direction rule such as
\texttt{q\_A\textgreater{}0.5} has expected success rate 50 percent
under this null. That fact is not a failure of calibration; it is the
signature of a boundary statistic.

The same argument extends to a linear boundary in a multivariate
Gaussian posterior. Let

\begin{Shaded}
\begin{Highlighting}[]
\NormalTok{x\_hat \textasciitilde{} Normal(theta0, Sigma\_data)}
\NormalTok{theta | x\_hat \textasciitilde{} Normal(x\_hat, Sigma\_post)}
\NormalTok{g(theta) = b\^{}T theta {-} c.}
\end{Highlighting}
\end{Shaded}

Then

\begin{Shaded}
\begin{Highlighting}[]
\NormalTok{q\_A = P(g(theta)\textgreater{}0 | x\_hat)}
\NormalTok{    = Phi((b\^{}T x\_hat {-} c) / sqrt(b\^{}T Sigma\_post b)).}
\end{Highlighting}
\end{Shaded}

If the truth lies on the boundary, \texttt{b\^{}T\ theta0=c}, and the
repeated-experiment variance in the boundary-normal direction matches
the posterior variance,
\texttt{b\^{}T\ Sigma\_data\ b\ =\ b\^{}T\ Sigma\_post\ b}, the same
Uniform(0,1) result follows.

This is the mathematical core of the paper. The observable class
probability is not wrong, but its interpretation changes. At a boundary,
it is a calibrated random variable whose null distribution must be
reported.

\subsection{3. One-Dimensional Boundary
Null}\label{one-dimensional-boundary-null}

The first simulation implements the one-dimensional model above with
\texttt{theta0=0}, \texttt{sigma\_data=1}, and \texttt{sigma\_post=1}.
The run uses 200000 repeated experiments. The baseline diagnostics are:

{\def\LTcaptype{none} % do not increment counter
\begin{longtable}[]{@{}ll@{}}
\toprule\noalign{}
Statistic & Value \\
\midrule\noalign{}
\endhead
\bottomrule\noalign{}
\endlastfoot
Mean of \texttt{q\_A} & 0.500960 \\
Variance of \texttt{q\_A} & 0.083559 \\
Uniform variance \texttt{1/12} & 0.083333 \\
Direction rate \texttt{q\_A\textgreater{}0.5} & 0.502015 \\
KS D vs Uniform(0,1) & 0.002498 \\
KS p vs Uniform(0,1) & 0.164593 \\
\end{longtable}
}

Figure 1 shows the corresponding histogram and empirical CDF. The
central point is not merely that the mean is near one half. The
important point is that the whole distribution is broad and
approximately uniform. A boundary-null class probability can take values
near 0.05, 0.5, or 0.95 across repeated experiments without implying
that the posterior calculation itself has failed.

The same simulation suite includes off-boundary and
posterior-width-mismatch checks. These are diagnostic controls. Moving
the truth away from the boundary introduces direction information;
making the posterior too wide or too tight distorts the
class-probability distribution. Those failures are useful because they
show that the uniform result is not a generic artifact. It depends on
the truth being at the boundary and on the calibration of
repeated-experiment and posterior uncertainty along the boundary-normal
direction.

\subsection{4. Cosmology-Shaped w0-wa
Analogue}\label{cosmology-shaped-w0-wa-analogue}

The second simulation translates the same idea into a two-dimensional
geometry resembling the local w0-wa plane. Define coordinates centered
on Lambda-CDM:

\begin{Shaded}
\begin{Highlighting}[]
\NormalTok{u = w0 + 1,}
\NormalTok{v = wa,}
\NormalTok{Lambda{-}CDM = (u,v) = (0,0).}
\end{Highlighting}
\end{Shaded}

The toy boundary score is

\begin{Shaded}
\begin{Highlighting}[]
\NormalTok{s(theta) = alpha u + v,}
\end{Highlighting}
\end{Shaded}

with \texttt{RIP-like} on the positive side and \texttt{DECAY-like} on
the negative side. This is not a physical CPL fate classifier. It is a
controlled boundary analogue in which the boundary-normal score is
explicit and Lambda-CDM lies on the boundary.

For 50000 repeated experiments, the baseline diagnostics are:

{\def\LTcaptype{none} % do not increment counter
\begin{longtable}[]{@{}ll@{}}
\toprule\noalign{}
Statistic & Value \\
\midrule\noalign{}
\endhead
\bottomrule\noalign{}
\endlastfoot
Mean of \texttt{P(RIP-like)} & 0.501283 \\
Variance & 0.083260 \\
Uniform variance \texttt{1/12} & 0.083333 \\
Direction rate \texttt{P(RIP-like)\textgreater{}0.5} & 0.505160 \\
KS D vs Uniform(0,1) & 0.005944 \\
KS p vs Uniform(0,1) & 0.058189 \\
\end{longtable}
}

The same Lambda-CDM truth can yield low, near-boundary, or high
posterior class masses depending on the noisy posterior center. the
two-dimensional experiment selected three examples from the same
repeated-experiment ensemble:

{\def\LTcaptype{none} % do not increment counter
\begin{longtable}[]{@{}llll@{}}
\toprule\noalign{}
Example & w0\_hat & wa\_hat & \texttt{P(RIP-like)} \\
\midrule\noalign{}
\endhead
\bottomrule\noalign{}
\endlastfoot
low RIP-like probability & -0.86259 & -0.45695 & 0.08001 \\
near boundary & -1.05691 & 0.06830 & 0.50000 \\
high RIP-like probability & -1.08011 & 0.38820 & 0.92001 \\
\end{longtable}
}

This result does not claim that the physical CPL fate boundary is
exactly linear or that the real universe is exactly on it. It
establishes the local statistical mechanism: once a posterior overlaps a
classification boundary, the boundary-normal projection controls the
class mass. A visually realistic w0-wa posterior can therefore yield
unstable fate labels even under a stable null.

\subsection{5. Fate Classifier as a Threshold
Map}\label{fate-classifier-as-a-threshold-map}

A physical fate classifier generally evaluates continuous
future-behavior diagnostics and then returns a discrete label. We
represent this by four margins:

\begin{Shaded}
\begin{Highlighting}[]
\NormalTok{crunch\_margin}
\NormalTok{rip\_margin}
\NormalTok{ds\_margin}
\NormalTok{decay\_margin}
\end{Highlighting}
\end{Shaded}

and a priority rule:

\begin{Shaded}
\begin{Highlighting}[]
\NormalTok{if crunch\_margin \textgreater{} 0:}
\NormalTok{    CRUNCH}
\NormalTok{elif rip\_margin \textgreater{} 0:}
\NormalTok{    RIP}
\NormalTok{elif ds\_margin \textgreater{} 0:}
\NormalTok{    DS}
\NormalTok{elif decay\_margin \textgreater{} 0:}
\NormalTok{    DECAY}
\NormalTok{else:}
\NormalTok{    OTHER}
\end{Highlighting}
\end{Shaded}

The labels retain their conceptual roles: \texttt{CRUNCH} for collapse,
\texttt{RIP} for a finite-time divergence, \texttt{DS} for de
Sitter-like asymptotics, \texttt{DECAY} for a dark-energy-decay or
eternal-expansion-compatible branch, and \texttt{OTHER} as an audit
channel. The method does not require the reader to accept one specific
physical implementation of these labels. It requires that whatever
classifier is used be treated as a threshold map with auditable
boundaries.

The the classifier-framework example example posterior produces mixed
class mass:

{\def\LTcaptype{none} % do not increment counter
\begin{longtable}[]{@{}ll@{}}
\toprule\noalign{}
Label & Posterior fraction \\
\midrule\noalign{}
\endhead
\bottomrule\noalign{}
\endlastfoot
CRUNCH & 0.00000 \\
RIP & 0.36857 \\
DS & 0.03286 \\
DECAY & 0.50729 \\
OTHER & 0.09129 \\
\end{longtable}
}

The associated path diagnostic shows a continuous descriptor path whose
label jumps from \texttt{DECAY} to \texttt{RIP} as the RIP margin
crosses zero. This is the classifier-level version of the
boundary-statistic problem. Smooth movement in descriptors can produce
discontinuous changes in labels, and posterior class masses are
integrals over those discontinuous regions.

\subsection{6. Finite-Mock Null
Calibration}\label{finite-mock-null-calibration}

Boundary calibration requires an explicit null distribution. This is
close in spirit to predictive calibration and posterior predictive
checking \citep{Dawid1984, Rubin1984, GelmanMengStern1996}. For a
predeclared classifier \texttt{C}, label \texttt{L}, and observed data
set, compute the raw class mass

\begin{Shaded}
\begin{Highlighting}[]
\NormalTok{p\_obs = P(C(theta)=L | observed data).}
\end{Highlighting}
\end{Shaded}

Then choose the relevant null model, generate \texttt{N} mock data sets,
and run the identical inference and classification pipeline on every
mock:

\begin{Shaded}
\begin{Highlighting}[]
\NormalTok{p\_j = P(C(theta)=L | mock\_j),  j=1,...,N.}
\end{Highlighting}
\end{Shaded}

The lower-tail, upper-tail, and two-sided finite-mock calibrated
p-values are:

\begin{Shaded}
\begin{Highlighting}[]
\NormalTok{p\_lower = (\#\{p\_j \textless{}= p\_obs\} + 1) / (N + 1)}
\NormalTok{p\_upper = (\#\{p\_j \textgreater{}= p\_obs\} + 1) / (N + 1)}
\NormalTok{p\_two\_sided = min(1, 2 * min(p\_lower, p\_upper)).}
\end{Highlighting}
\end{Shaded}

The plus-one rule prevents zero p-values from finite mock ensembles,
matching the finite randomization principle that simulation p-values
should not be reported as zero \citep{PhipsonSmyth2010}. The minimum
reportable one-sided tail probability is

\begin{Shaded}
\begin{Highlighting}[]
\NormalTok{p\_floor = 1 / (N + 1).}
\end{Highlighting}
\end{Shaded}

This calibration does not replace the raw class mass. It answers a
different question. The raw mass says how much posterior probability
lies in a label region. The calibrated tail depth says how unusual that
raw mass is under the declared null.

The synthetic the finite-mock demonstration demonstration used 100
finite null mocks and three observed values. The results were:

{\def\LTcaptype{none} % do not increment counter
\begin{longtable}[]{@{}lllll@{}}
\toprule\noalign{}
Case & Raw probability & Lower-tail p & Upper-tail p & Two-sided p \\
\midrule\noalign{}
\endhead
\bottomrule\noalign{}
\endlastfoot
Central boundary-like & 0.500 & 0.44554 & 0.56436 & 0.89109 \\
Low-tail example & 0.018 & 0.03960 & 0.97030 & 0.07921 \\
High-tail example & 0.982 & 0.98020 & 0.02970 & 0.05941 \\
\end{longtable}
}

The synthetic examples are not cosmology claims. They illustrate the
reporting grammar that the real-pipeline case study then uses.

\subsection{7. Real-Pipeline LCDM Mock Case
Study}\label{real-pipeline-lcdm-mock-case-study}

The case study applies the calibration protocol to existing outputs from
a cosmology pipeline. It uses 100 noisy LCDM mocks as the null
distribution and keeps the zero-noise mock separate as a mechanism
diagnostic. The analyzed quantities are

\begin{Shaded}
\begin{Highlighting}[]
\NormalTok{P\_heat = P(DS) + P(DECAY),}
\NormalTok{P\_RIP  = P(RIP).}
\end{Highlighting}
\end{Shaded}

The noisy mocks are the finite null distribution. The zero-noise mock is
not included in the null tail counts because it answers a different
question: what the deterministic pipeline does when the mock data
exactly match the null without noise.

The case-study results are:

{\def\LTcaptype{none} % do not increment counter
\begin{longtable}[]{@{}ll@{}}
\toprule\noalign{}
Quantity & Value \\
\midrule\noalign{}
\endhead
\bottomrule\noalign{}
\endlastfoot
Noisy LCDM mocks & 100 \\
Null \texttt{P\_heat} mean & 0.52515 \\
Null \texttt{P\_heat} standard error & 0.03026 \\
Null \texttt{P\_heat} KS p vs Uniform & 0.25266 \\
Direction rate \texttt{P\_heat\ \textgreater{}\ 0.5} & 0.50000 \\
Observed \texttt{P\_heat} & 0.99821 \\
Observed \texttt{P\_RIP} & 0.001792 \\
Noisy mocks with \texttt{P\_RIP\ \textless{}=\ observed} & 0 / 100 \\
Plus-one lower-tail p for \texttt{P\_RIP} & 0.00990 \\
Zero-count 95\% upper bound & 0.02951 \\
Zero-noise mock000 \texttt{P\_heat} & 0.52275 \\
Zero-noise mock000 \texttt{P\_RIP} & 0.47577 \\
Max OTHER fraction in noisy mocks & 0.001416 \\
Max boundary fraction in noisy mocks & 0.000871 \\
\end{longtable}
}

The direction layer is deliberately uninformative under this LCDM
boundary null: \texttt{P\_heat\textgreater{}0.5} in 50/100 noisy mocks,
and the zero-noise mock sits near a 52/48 heat/RIP split. This would
look poor if judged by a high fixed direction accuracy threshold, but it
is the expected boundary behavior.

The depth layer says something different. The observed \texttt{P\_RIP}
is lower than all 100 noisy LCDM mocks. With plus-one smoothing this
becomes a one-sided lower-tail p-value of 0.00990, not zero. The
zero-count 95 percent upper bound is 0.02951. This is an interpretable
tail-depth statement under the declared null and finite mock count.

The correct reading is therefore narrow:

\begin{Shaded}
\begin{Highlighting}[]
\NormalTok{Under the declared LCDM null and this pipeline, the direction layer is unstable,}
\NormalTok{while the observed P\_RIP is unusually low relative to the 100 noisy null mocks.}
\end{Highlighting}
\end{Shaded}

It is not a direct forecast of the real universe's fate. It is not, by
itself, evidence that a larger dark-energy model is preferred. Those
questions belong to the Evidence and Horizon layers described next.

\subsection{8. Reporting Standard}\label{reporting-standard}

The reporting unit for boundary-sensitive fate probabilities should have
five layers:

{\def\LTcaptype{none} % do not increment counter
\begin{longtable}[]{@{}
  >{\raggedright\arraybackslash}p{(\linewidth - 6\tabcolsep) * \real{0.2500}}
  >{\raggedright\arraybackslash}p{(\linewidth - 6\tabcolsep) * \real{0.2500}}
  >{\raggedright\arraybackslash}p{(\linewidth - 6\tabcolsep) * \real{0.2500}}
  >{\raggedright\arraybackslash}p{(\linewidth - 6\tabcolsep) * \real{0.2500}}@{}}
\toprule\noalign{}
\begin{minipage}[b]{\linewidth}\raggedright
Layer
\end{minipage} & \begin{minipage}[b]{\linewidth}\raggedright
Question
\end{minipage} & \begin{minipage}[b]{\linewidth}\raggedright
Required report
\end{minipage} & \begin{minipage}[b]{\linewidth}\raggedright
Failure mode
\end{minipage} \\
\midrule\noalign{}
\endhead
\bottomrule\noalign{}
\endlastfoot
Direction & Which label has more posterior mass? & raw P(label
\textbar{} data); direction threshold, usually 0.5; same quantity under
null mocks & Treating P \textgreater{} 0.5 as a discovery near a
boundary. \\
Depth & How unusual is the raw probability under the null? & predeclared
null model; mock count and finite-mock floor; lower-tail, upper-tail,
and two-sided calibrated p-values & Reporting a raw tail probability
without null calibration. \\
Evidence & Does the enlarged model deserve preference? &
likelihood-ratio or delta chi2 where appropriate; AIC/BIC or Bayes
factor where available; clear separation from fate-class probability &
Confusing a class probability with evidence for the model. \\
Boundary & Is the label sensitive to thresholds or classifier defects? &
boundary-flip fraction or threshold sensitivity; OTHER fraction and
audit rule; known classifier limitations & Hiding that posterior mass
lies near a label boundary. \\
Horizon & Is the fate decision inside the data-constrained region? &
future scale used by the classifier; constraint horizon or equivalent
diagnostic; statement when the horizon diagnostic is unavailable &
Interpreting unconstrained far-future extrapolation as data. \\
\end{longtable}
}

The minimum prose should contain three distinct statements:

\begin{Shaded}
\begin{Highlighting}[]
\NormalTok{Raw class mass:}
\NormalTok{The posterior mass in label L is ...}

\NormalTok{Calibrated boundary statistic:}
\NormalTok{Under the predeclared null, the observed class probability lies in the ...}
\NormalTok{tail with finite{-}mock{-}calibrated p = ...}

\NormalTok{Guardrail:}
\NormalTok{This is not, by itself, a direct forecast of the universe\textquotesingle{}s fate.}
\end{Highlighting}
\end{Shaded}

The the real-pipeline case study case study fills Direction, Depth, and
Boundary. It does not fill Evidence and Horizon. That incompleteness is
intentional. A method paper should not silently promote an available
class-mass result into an unavailable model preference or
extrapolation-horizon result.

For the case study, the layer fill is:

{\def\LTcaptype{none} % do not increment counter
\begin{longtable}[]{@{}
  >{\raggedright\arraybackslash}p{(\linewidth - 2\tabcolsep) * \real{0.5000}}
  >{\raggedright\arraybackslash}p{(\linewidth - 2\tabcolsep) * \real{0.5000}}@{}}
\toprule\noalign{}
\begin{minipage}[b]{\linewidth}\raggedright
Layer
\end{minipage} & \begin{minipage}[b]{\linewidth}\raggedright
Case-study fill
\end{minipage} \\
\midrule\noalign{}
\endhead
\bottomrule\noalign{}
\endlastfoot
Direction & Observed P\_heat=0.99821; LCDM null has
P\_heat\textgreater0.5 in 50/100 mocks. \\
Depth & Observed P\_RIP=0.00179; 0/100 null mocks are lower; plus-one
p=0.00990, 95\% CP upper=0.02951. \\
Evidence & Not evaluated by the real-pipeline case study figure; a
science paper must add likelihood-ratio, AIC/BIC, or Bayes-factor
evidence separately. \\
Boundary & max OTHER=0.00142; max boundary fraction=0.00087; mock000
P\_heat=0.52275. \\
Horizon & Not evaluated by the real-pipeline case study figure; a
science paper must report constraint horizon or explicitly mark it
unavailable. \\
\end{longtable}
}

This is the practical contribution of the paper. The standard provides a
citable format for future fate analyses: raw class mass, null-calibrated
depth, model evidence, boundary sensitivity, and horizon validity must
be separately visible.

\subsection{9. Robustness and Diagnostic
Failures}\label{robustness-and-diagnostic-failures}

A method that only shows a favorable toy example is not enough. the
robustness exercise therefore asks which perturbations preserve the
boundary-null behavior and which correctly fail diagnostics.

{\def\LTcaptype{none} % do not increment counter
\begin{longtable}[]{@{}
  >{\raggedright\arraybackslash}p{(\linewidth - 10\tabcolsep) * \real{0.1667}}
  >{\raggedright\arraybackslash}p{(\linewidth - 10\tabcolsep) * \real{0.1667}}
  >{\raggedright\arraybackslash}p{(\linewidth - 10\tabcolsep) * \real{0.1667}}
  >{\raggedright\arraybackslash}p{(\linewidth - 10\tabcolsep) * \real{0.1667}}
  >{\raggedright\arraybackslash}p{(\linewidth - 10\tabcolsep) * \real{0.1667}}
  >{\raggedright\arraybackslash}p{(\linewidth - 10\tabcolsep) * \real{0.1667}}@{}}
\toprule\noalign{}
\begin{minipage}[b]{\linewidth}\raggedright
Variant
\end{minipage} & \begin{minipage}[b]{\linewidth}\raggedright
Mean
\end{minipage} & \begin{minipage}[b]{\linewidth}\raggedright
Variance
\end{minipage} & \begin{minipage}[b]{\linewidth}\raggedright
Direction rate
\end{minipage} & \begin{minipage}[b]{\linewidth}\raggedright
KS p
\end{minipage} & \begin{minipage}[b]{\linewidth}\raggedright
Status
\end{minipage} \\
\midrule\noalign{}
\endhead
\bottomrule\noalign{}
\endlastfoot
calibrated & 0.498806 & 0.083597 & 0.497925 & 0.147 & pass \\
shifted & 0.499364 & 0.083356 & 0.497250 & 0.245 & pass \\
off-boundary & 0.638641 & 0.075633 & 0.691025 & 0 & diagnostic
failure \\
too wide & 0.499690 & 0.049937 & 0.499387 & 0 & diagnostic failure \\
too tight & 0.499530 & 0.117285 & 0.498800 & 0 & diagnostic failure \\
curved beta=1.5 & 0.524333 & 0.086348 & 0.540333 & 2.51e-06 & curvature
diagnostic \\
\end{longtable}
}

The calibrated boundary and shifted-boundary cases pass. This shows that
the uniform result is not tied to the arbitrary location of the
threshold; if the truth and the threshold move together and calibration
is correct, the boundary-null result remains.

The off-boundary case fails, as it should. Moving the truth away from
the threshold creates direction information. The posterior-width
mismatch cases also fail. A posterior that is too wide becomes
under-confident and has too little variance in class probability; a
posterior that is too tight becomes over-confident and has too much
variance. These failures are diagnostic features, not method failures.

The curved-boundary case is labeled as a curvature diagnostic. Its mean
is 0.524 and its KS p-value is 2.51e-06. The method does not claim that
arbitrary curvature preserves uniformity. It says that curvature,
posterior calibration, and distance from the boundary must be reported
or tested instead of hidden.

Finite mock resolution is another robustness constraint. For an observed
raw probability around 0.018, 100 mocks have a plus-one p-value floor of
1/101 = 0.00990. The probability of seeing zero mocks below that
observed value under a Uniform null is still about 0.163 at
\texttt{N=100}; it falls only with larger mock ensembles. This means
finite null ensembles should be reported with their resolution, not
treated as unlimited tail probes.

\subsection{10. Claim Ledger and Proof
Boundary}\label{claim-ledger-and-proof-boundary}

The paper's claims are deliberately limited:

{\def\LTcaptype{none} % do not increment counter
\begin{longtable}[]{@{}
  >{\raggedright\arraybackslash}p{(\linewidth - 4\tabcolsep) * \real{0.3333}}
  >{\raggedright\arraybackslash}p{(\linewidth - 4\tabcolsep) * \real{0.3333}}
  >{\raggedright\arraybackslash}p{(\linewidth - 4\tabcolsep) * \real{0.3333}}@{}}
\toprule\noalign{}
\begin{minipage}[b]{\linewidth}\raggedright
Status
\end{minipage} & \begin{minipage}[b]{\linewidth}\raggedright
Claim
\end{minipage} & \begin{minipage}[b]{\linewidth}\raggedright
Support or boundary
\end{minipage} \\
\midrule\noalign{}
\endhead
\bottomrule\noalign{}
\endlastfoot
supported & Near a calibrated critical boundary, posterior class
probabilities can be broad and approximately uniform under repeated
experiments. & the one-dimensional experiment mean=0.50096,
variance=0.08356, KS p=0.16459; the two-dimensional experiment KS
p=0.05819. \\
supported & A 50 percent direction rate is not automatically a failed
classifier near the boundary. & the real-pipeline case study LCDM mocks
have P\_heat\textgreater0.5 in 0.50 of mocks while the P\_heat null
shape passes KS p=0.25266. \\
supported & Tail depth can still carry information when direction is
unstable. & Observed P\_RIP=0.00179 is below 0/100 noisy LCDM mocks;
plus-one p=0.00990. \\
supported & The reporting unit must separate Direction, Depth, Evidence,
Boundary, and Horizon. & the reporting-standard exercise defines all
five layers and marks Evidence and Horizon as unfilled in the case
study. \\
not claimed & The method establishes the real universe's final fate. &
Excluded by the problem-definition section non-goals and by the
real-pipeline case study scope. \\
not claimed & The case study proves dynamical dark energy or model
preference. & Evidence layer is explicitly unavailable in the
reporting-standard exercise and must be supplied by a separate science
analysis. \\
not claimed & Curved or miscalibrated boundaries always preserve
uniformity. & the robustness exercise shows off-boundary truth and
posterior-width mismatch fail; beta=1.5 curvature is labeled a
diagnostic. \\
\end{longtable}
}

The negative entries are as important as the positive ones. This method
paper does not establish the real universe's final fate. It does not
prove dynamical dark energy or any model preference. It does not claim
that curved or miscalibrated boundaries always preserve uniformity. It
provides a calibration and reporting framework for a specific
statistical failure mode: interpreting boundary-sensitive class masses
as direct fate forecasts.

\subsection{11. Limitations}\label{limitations}

The main limitation is that the toy proofs are local. The
one-dimensional and linear w0-wa analogues explain the boundary-normal
mechanism, but a physical fate classifier may have curved, intersecting,
or prior-truncated boundaries. The response is not to discard the
method; it is to report the relevant boundary diagnostics and to treat
strong curvature as a diagnostic condition.

The second limitation is finite mock resolution. With 100 mocks, the
smallest reportable one-sided p-value is 0.00990. This is adequate for
demonstrating the reporting protocol, but not for resolving much deeper
tails. Larger mock ensembles would be needed for stronger tail claims.

The third limitation is the case-study scope. The LCDM mock case
demonstrates direction-depth separation for one pipeline and one
declared null. It does not establish that all cosmological pipelines,
data combinations, or classifiers would produce the same calibrated tail
depth.

The fourth limitation is the unfilled Evidence and Horizon layers.
Evidence requires likelihood-ratio, information-criterion, Bayes-factor,
or equivalent model-comparison work. Horizon requires an assessment of
whether the future scale used by the fate classifier lies inside a
data-constrained regime. The case study intentionally does not fill
those rows.

Finally, this paper does not validate compressed CMB likelihoods as
equivalent to full Planck analyses, nor does it settle the current
observational debate about evolving dark energy. Those are science-paper
tasks. The present method paper supplies the reporting discipline needed
before such results are turned into fate statements.

\subsection{12. Conclusion}\label{conclusion}

Posterior fate probabilities near a critical boundary should be treated
as boundary statistics. In the calibrated boundary null, broad class
probabilities and 50 percent direction rates are expected. That does not
make the raw class mass useless; it changes the reporting question. The
raw mass must be paired with a null distribution, finite-mock tail
calibration, boundary diagnostics, model evidence, and an
extrapolation-horizon statement.

The proposed five-layer standard gives a compact rule: report Direction,
Depth, Evidence, Boundary, and Horizon separately. The LCDM mock case
study shows why this matters. Direction can be unstable under the null,
while an observed tail depth can still be calibrated and reported.
Keeping those statements separate prevents both false pessimism about a
working classifier and false confidence about the universe's final fate.

\subsection{Figure Captions and Source
Files}\label{figure-captions-and-source-files}

{\def\LTcaptype{none} % do not increment counter
\begin{longtable}[]{@{}
  >{\raggedright\arraybackslash}p{(\linewidth - 4\tabcolsep) * \real{0.3333}}
  >{\raggedright\arraybackslash}p{(\linewidth - 4\tabcolsep) * \real{0.3333}}
  >{\raggedright\arraybackslash}p{(\linewidth - 4\tabcolsep) * \real{0.3333}}@{}}
\toprule\noalign{}
\begin{minipage}[b]{\linewidth}\raggedright
Figure
\end{minipage} & \begin{minipage}[b]{\linewidth}\raggedright
Caption target
\end{minipage} & \begin{minipage}[b]{\linewidth}\raggedright
Source
\end{minipage} \\
\midrule\noalign{}
\endhead
\bottomrule\noalign{}
\endlastfoot
1 & One-dimensional boundary-null calibration &
figures/phase1\_toy\_boundary\_uniformity.png \\
2 & Two-dimensional w0-wa boundary analogue &
figures/phase2\_w0wa\_boundary.png \\
3 & Abstract fate-classifier framework &
figures/phase3\_classifier\_framework.png \\
4 & Finite-mock null-calibration protocol &
figures/phase4\_null\_calibration.png \\
5 & Real-pipeline LCDM mock case study &
figures/phase5\_cosmology\_mock\_case.png \\
6 & Five-layer reporting standard &
figures/phase6\_reporting\_standard.png \\
A1 & Robustness diagnostics & figures/phase7\_robustness.png \\
\end{longtable}
}

\subsection{Appendix A. Robustness
Details}\label{appendix-a.-robustness-details}

The one-dimensional robustness checks use the same Gaussian boundary
statistic as Section 3. The pass rule requires the mean to be within
0.02 of 0.5, the variance to be within 0.02 of \texttt{1/12}, the
direction rate to be within 0.02 of 0.5, and the KS p-value against
Uniform(0,1) to be at least 0.01. The calibrated and shifted-boundary
cases pass these checks. The off-boundary, posterior-too-wide, and
posterior-too-tight cases fail in the intended directions.

The curved-boundary experiment uses a two-dimensional score

\begin{Shaded}
\begin{Highlighting}[]
\NormalTok{score = v + alpha*u + beta*u\^{}2}
\end{Highlighting}
\end{Shaded}

with truth at \texttt{(u,v)=(0,0)} and \texttt{beta=1.5}. This setting
is intentionally strong enough to trigger a curvature diagnostic. A
physical application should therefore report local boundary curvature or
use mocks generated through the full classifier rather than relying only
on a linear approximation.

\subsection{Appendix B. Reporting
Checklist}\label{appendix-b.-reporting-checklist}

Use this checklist whenever a fate probability is reported near a
critical boundary:

\begin{enumerate}
\def\labelenumi{\arabic{enumi}.}
\tightlist
\item
  State the raw posterior class mass and the classifier used.
\item
  State the null model and why it is the relevant boundary or
  near-boundary null.
\item
  State the number of mocks and the finite-mock p-value floor.
\item
  Report lower-tail, upper-tail, and two-sided calibrated p-values.
\item
  Report direction behavior under the null separately from tail depth.
\item
  Report model evidence separately from class mass.
\item
  Report boundary sensitivity, OTHER fraction, or equivalent classifier
  audit.
\item
  Report the extrapolation horizon or mark it unavailable.
\item
  Include the guardrail: the class mass is not by itself a direct
  final-fate forecast.
\end{enumerate}

\subsection{References}\label{references}

The final bibliography is supplied in \texttt{refs\_method.bib} for the
submission package.

\bibliography{/Users/nocide/cosmo_fate/paper/method/submission_final/refs.bib}

\end{document}
